\section{Limitations and Future Work}
\label{sec:limitations}

\paragraph{The benefit of user-guided pivots is not yet causally evaluated.}
The retained traces show evidence-backed changes to operational objectives,
constraints, and verification obligations, but they do not include a prospective
human-subject comparison. We therefore do not claim that the current experiments
prove improved expert understanding. A direct evaluation should compare
\systemname with fixed-prompt agents and ordinary interactive assistants on expert
calibration, time to identify the binding constraint, quality of contract revisions,
number of repeated misconceptions, and final task value.
Likewise, internal-expert observations about appropriate stopping motivate the
analysis in Section~\ref{sec:discussion} but are not reported as a satisfaction rate
or comparative user-study result.

\paragraph{Contract refinement can still fail.}
Evidence, provenance, and authority reduce goal drift but do not eliminate it. An
expert may approve a poorly framed tradeoff, Reviewer may miss a contradiction, or
the recorded standing intent may omit a tacit requirement. Future work should test
adversarial contract changes, invariant-preservation checks, disagreement handling,
and rollback of harmful refinements. High-impact domains also require explicit
human sign-off policies rather than relying on Manager judgment alone.

The task-contract notation is an analytical abstraction over several persisted
runtime surfaces, not a claim that the implementation currently exposes one atomic
$K_t$ object or a universal automated user-admission API. A future implementation
should unify those fields and make authorization and rollback directly queryable.

\paragraph{Verification is only as sound as its evidence boundary.}
An executable test, formal checker, benchmark, or model Reviewer can all encode the
wrong property. The framework makes verifier revision possible but does not confer
correctness automatically. System-verification studies should separately seed bugs
in code, bugs in specifications, and mismatches between them, then measure whether
the runtime localizes the faulty side instead of merely making the current checker
pass.

\paragraph{Attribution and task sequence.}
The SWE-Bench Pro experiment compares the complete \systemname runtime with Direct
Copilot. Reviewer routing is adaptive, the 22 Waves follow one task order, and
per-Wave Direct-Copilot Token/time records are unavailable. The current results
therefore characterize whole-system behavior over this sequence. Matched
frozen-state runs, randomized task orders, and randomized review routing will
separate the contributions of persistent state, role separation, and review.

\paragraph{Generality across models, domains, and evaluators.}
The seven benchmark arenas use different metrics, hardware, and execution
protocols, while the mathematical analysis follows one vertical campaign. Results
are interpreted in their native units rather than combined into a universal score.
Repeated campaigns, additional backbones, and external domain review will measure
how the runtime transfers across task families and research standards.

\paragraph{Metric calibration and model transfer.}
The efficiency analysis labels whole role calls by mission purpose; finer segment
labels would sharpen the reasoning, action, and verification factors. Reviewer
sensitivity and false-acceptance rate require randomized routing with an external
verifier, and the counterfactual reuse value $G_L$ requires paired future tasks with
and without a state update. The runtime-to-model path is currently a training
direction: held-out trajectory splits and scaffold-removal studies will measure how
much of the long-horizon control loop can be internalized by the model.

\paragraph{Paper-production trajectories are observational.}
The six paper projects were selected from one shared research environment and do
not include a matched human-authoring baseline. Campaign-hours overlap in calendar
time, stored review snapshots are model-generated rather than independent human
reviews, and runtime logging evolved across projects. The case study therefore
supports end-to-end lifecycle and recovery claims, not paper acceptance, scientific
novelty, or superiority to human research teams. Recorded costs in the public data
also exclude utility calls without a priced event.

\paragraph{Artifact evidence levels.}
The FLA case is linked to a public merged pull request and archived timing and
correctness evidence. The ACE-1 description and 75/80 score come from the supplied
project presentation; the report bundle does not include enough raw PPA and
throughput material for independent recomputation. ACE-1 is therefore evidence of
an artifact-producing chip-design workflow, not a reproduced SOTA comparison.
